% !TeX spellcheck = de_DE
\addscenariosection{1}{Kooperatives Szenario}{Zorn des bösen Zauberers}{\images/zorn_des_b_sen_zaubererers.png}

\begin{multicols*}{2}

\textbf{Autor:} Dwarftoe

\textit{Unzählige Jahre lang hatte ein böser Zauberer in seinem bedrohlichen Turm auf einer verlassenen Insel gelebt, weit vor den Küsten einiger benachbarter Königreiche. Von den Küstendörfern aus blickten die Menschen oft zum Horizont und sahen eine gewaltige schwarze Wolke über der Insel hängen, durchzuckt von Blitzen. Aus diesem verfluchten Ort kamen eisige Winde, die die Kälte des Todes selbst mit sich zu tragen schienen. In ihrer Angst wandten sich viele an einen uralten Meeresgott. Sein vergessener Tempel, halb von den Wellen verschlungen, stand nicht weit von der Küste entfernt; seine verwitterten Steine ragten aus dem dunklen Wasser wie die Ruinen eines längst vergangenen Zeitalters. Die Menschen beteten, dass der alte Gott sie vor dem Zorn des Zauberers schützen möge. Doch kein Schiff, das es wagte, zur Insel zu segeln, kehrte jemals zurück. Jedes Schiff, das ausgesandt wurde, um sich dem Zauberer entgegenzustellen, fand ein schreckliches Ende. Manche wurden von Blitzen verschlungen, die aus dem sturmumtosten Himmel herabfuhren, während andere von monströsen Wellen erfasst und gegen die zerklüfteten Klippen der Insel geschleudert wurden. Doch dann änderte sich ohne jede Vorwarnung alles. Die Winde begannen mit unnatürlicher Wut zu heulen, und die Kälte wurde mit jeder Stunde schlimmer. Ein gewaltiger Schneesturm senkte sich über die Königreiche und fegte wie eine Armee weißer Schatten über das Land. Felder verschwanden unter Bergen von Schnee. Flüsse froren vollständig zu. Wälder wurden zu stillen Königreichen aus Eis. Selbst das Meer begann zu gefrieren. Die Menschen an der Küste hatten keine andere Wahl, als in ihre Häuser zu fliehen und die Türen gegen den heulenden Sturm zu verbarrikadieren. Wer draußen überrascht wurde, ging bald in der erbarmungslosen Kälte verloren. Und doch blieb ein Ort unberührt. Rund um den uralten Tempel des Meeresgottes schien der Sturm nachzulassen. Das Wasser wollte nicht vollständig gefrieren, und die Winde wurden schwächer, wenn sie über die heiligen Ruinen hinwegzogen. In der Nähe lag eine kleine, einsame Insel, auf der seit vielen Jahren eine alte Hexe zurückgezogen lebte. Auf unerklärliche Weise hielten die Hexe und der uralte Meeresgott gemeinsam den Schneesturm in Schach und schützten die dahinterliegenden Länder vor seinem eisigen Griff. Doch nun schwand ihre Kraft. Die Winde begannen, die Schutzbarriere zu durchbrechen. Das Eis kroch immer weiter über das Land und verschlang nach und nach Wälder, Felder und Dörfer. Die uralten Schutzkräfte, die Jahrhunderten getrotzt hatten, wurden schwächer, und bald könnten selbst der Tempel und die Insel der Hexe vom Frost verschlungen werden. Die Herrscher der Königreiche standen nun vor einer schrecklichen Entscheidung. Es blieb nur wenig Zeit. Sie mussten ihre Rivalitäten beiseitelegen, ihre Armeen sammeln und in die gefrorenen Einöden aufbrechen, um den bösen Zauberer an der Quelle seiner Macht zu stellen. Denn wenn sie scheiterten, würde der Schneesturm alles verschlingen. Königreiche würden unter dem Eis verschwinden. Uralte Magie würde verstummen. Und alle, die diese Länder ihre Heimat nannten, würden einem endlosen Winter zum Opfer fallen.}

\subsection*{\MakeUppercase{Szenariodauer}}

15 Runden

\subsection*{\MakeUppercase{Spieleraufstellung}}

\textbf{Spieleranzahl:} 2--3

\textbf{Startressourcen:} 13 \svg{gold}, 4 \svg{building_materials}, 1 \svg{valuables}

\textbf{Start-Einkommen:} 10 \svg{gold}, 0 \svg{building_materials}, 0 \svg{valuables}

\textbf{Starteinheiten:}
\begin{itemize}
  \item Ein ``Rudel'' der günstigsten \svg{bronze}-Einheiten
  \item Wenige der teuersten \svg{bronze}-Einheiten
\end{itemize}

\textbf{Stadtgebäude:} \svg{bronze}~Gebäude, Rathaus, \svg[12]{building_special_gears}~Gebäude

\textbf{Gebietsvorrat:} Jeder Spieler nimmt 1 zufälliges fernes Gebiet (II--III). Legt alle 4 fernen Gebiete (II--III) aus der Turm-Erweiterung mit Eis-Hintergrund für eine spätere Verwendung beiseite.

\subsection*{\MakeUppercase{Kartenaufbau}}

Nehmt die folgenden Gebiete und ordnet sie wie in der Szenariokarte dargestellt an ($P$ steht für die Anzahl der Spieler):
\begin{itemize}
  \item $P$ × Startgebiete (I)
  \item 1 × Startgebiet (I) Turm
  \item $2P$ × ferne Gebiete (II--III)
  \item 2 × nahe Gebiete (IV--V) mit Eis-Hintergrund
  \item 1 × nahes Gebiet (IV--V) mit Elementarwind-Hintergrund (zählt als Eis)
  \item 1 × nahes Wassergebiet (IV--V) mit einer Hexenhütte in der Mitte
  \item 1 × zentrales Wassergebiet (VI--VII) mit dem Tempel des Meeresgottes
  \item 1 × zentrales Gebiet (VI--VII) mit Turmsiedlung
\end{itemize}

\subsection*{\MakeUppercase{Siegbedingungen}}
Besiegt den bösen Zauberer und erobert seine Stadt.

\subsection*{\MakeUppercase{Niederlagebedingungen}}
Am Ende von Runde 15 ist alles verloren.

\subsection*{\MakeUppercase{Zeitgesteuerte Ereignisse}}

\textbf{3. Runde:}
\begin{itemize}
  \item Tauscht ein zufälliges Gebiet II--III gegen ein zufälliges Eisgebiet (Gebiet II--III aus der Turm-Erweiterung). Nummeriert beispielsweise die Gebiete von 1 bis 6 und würfelt. Würfelt ihr eine 3, wird das dritte Gebiet ausgetauscht. Alle Würfel auf dem Gebiet werden entfernt und das Einkommen wird angepasst.
\end{itemize}

\textbf{6. Runde:}
\begin{itemize}
  \item Wie in Runde 3.
\end{itemize}

\textbf{9. Runde:}
\begin{itemize}
  \item Wie in Runde 3.
\end{itemize}

\textbf{12. Runde:}
\begin{itemize}
  \item Wie in Runde 3.
\end{itemize}

\subsection*{\MakeUppercase{Zusätzliche Regeln}}

\begin{itemize}
    \item Alle Spieler können einmal pro Runde Gold und Ressourcen tauschen. Einheiten, Artefakte und Zauber können ebenfalls getauscht werden, wenn sich ihre Helden auf benachbarten Feldern befinden.
    \item Kämpft ihr auf einem Eisgebiet, erleiden alle eure Einheiten zu Beginn 1 Schadenspunkt durch die Kälte.
    \item Befindet sich euer Held auf einem fernen Gebiet (II--III), während dieses in ein Eisgebiet umgewandelt wird, werdet ihr vom Schneesturm getroffen und eure schwächste Einheit stirbt.
    \item Der erste Held, der den Tempel des Meeresgottes besucht, erhält dessen Segen. Seine Einheiten erleiden bei Kämpfen auf Eisgebieten keinen anfänglichen Schadenspunkt mehr.
    \item Der erste Held, der die Insel mit der Hexenhütte besucht, erhält ihre Hilfe: eine fortgeschrittene Fähigkeitenkarte anstelle einer normalen sowie 3 Aufwertungs-Kronenmarker. Ihr könnt sie selbst behalten oder mit euren Kampfgefährten teilen.
    \item Bevor ihr die Stadt des Zauberers betreten könnt, müsst ihr seine Garnison besiegen (die Stadt auf dem Gebiet VI--VII neben dem Turm-Startgebiet). Die Garnison besteht aus 2 \svg{azure} Titanen, 1 \svg{golden} Naga, 1 \svg{silver} Dschinn und 1 \svg{silver} Magier (alles neutrale Einheiten) sowie Stadtmauern.
    \item Der Spieler, der die Garnison besiegt hat, kann zwischen zwei Optionen wählen: (1) Eine seiner \svg{bronze}- oder \svg{silver}-Einheiten kostenlos rekrutieren oder aufwerten oder eine \svg{golden}-Einheit für die Hälfte der Kosten rekrutieren oder aufwerten. (2) Den \svg{silver} Dschinn oder den \svg{silver} Magier kostenlos rekrutieren oder einen der \svg{azure} Titanen oder die \svg{golden} Naga für die Hälfte der Kosten rekrutieren.
    \item Wenn ihr gegen den Zauberer kämpft, müsst ihr alle 7 Einheiten der Turm-Fraktion auf der Rudel-Seite besiegen. Ihr beginnt mit zufällig ausgewählten 5 Einheiten; wird eine Einheit besiegt, rückt eine weitere in den Kampf nach. Die Stadt verfügt über Mauern und einen Geschützturm. Der Zauberer wird durch einen KI-Bot von Solmyr dargestellt. Ihr benötigt 8 Bot-Karten (1 Angriff, 1 Verteidigung, 2 Fähigkeiten, 4 Zauber). Die Fähigkeit ist Zauberei. Die Zauber sind: 1 × Kettenblitz, 1 × Blitz, 1 × Frostring, 1 × Magischer Pfeil.
    \item Verliert ein Spieler den Kampf, bleiben die besiegten Einheiten besiegt (ebenso Mauern und Geschützturm), aber die Bot-Karten werden für den nächsten Kampf neu gemischt.
\end{itemize}

\end{multicols*}

\vspace{3em}
\begin{center}
 \includegraphics[width=0.4\paperwidth]{\maps/zorn_des_b_sen_zaubererers_2p.png}
  \captionof{figure}{\textbf{SZENARIO FÜR 2 SPIELER}}
  \vspace{3em}
  \includegraphics[width=0.55\paperwidth]{\maps/zorn_des_b_sen_zaubererers_3p.png}
  \captionof{figure}{\textbf{SZENARIO FÜR 3 SPIELER}}
  \end{center}
